% ! TEX root = ../../main.tex

\subsection{Discussion}\label{sec:orca:design-discussion}


We now revisit ORCA's design to assess its guarantees, costs, and tradeoffs.

% ORCA's design is shaped by two considerations that warrant direct treatment: the consistency
% guarantees it provides, and the sources of perturbation it introduces. Both follow from the same
% premise --- BSP's synchronization structure makes certain properties cheap that would be expensive
% elsewhere, and ORCA's role is to exploit this structure rather than work against it.

% - Why guarantees? - ORCA's design choices make it unique. Not quite a database. It is a
% programmable network for transformation of telemetyr. Uplevel DAQ point (closest analogs are data
% acquisition pipelines --- frontends that discard data in FPGAs) Exploits homogeneity of a BSP
% application. All ranks are identical. Planning is easy. Starting point for exploiting in-network
% resources such as DPUs and programmable networks for an end-to-end usecase. BSP intrinsically has
% guarantees. Replication. Causal axes in async workloads are hard to define apriori. But BSP is
% delightfully simple. There is a single global logical clock. Understanding this clock is necessary
% for causal analyses and in-situ discarding. For example, saying that x mpi waits exceeded 10ms
% requires having all the data. Causal inference requires having all data: you cannot conclusively
% say which rank straggled without having all the data for a collective.

\pplushead{Rationale for Design Guarantees}
ORCA is best understood as a programmable network for BSP-aware in-situ trace analytics, not as
a conventional database for time-series data~\cite{:Prometheus,Agela2014:LDMS}.
While databases typically persist before analytics, ORCA enables in-situ workflows over a
tree-based overlay before ingesting processed and reduced output.
The closest analogs are \emph{Data Acquisition} (DAQ) pipelines in particle
accelerators~\cite{Aboli2016:DAQ-ATLAS,Bauer2012:DAQ-CMS}, where
FPGAs discard the bulk of sensor data at source, retaining only high-signal telemetry.
Its data and control plane guarantees are directly informed by observability requirements
given BSP execution.

\begin{itemize}[leftmargin=*]
  \item On the data path, the goal is to enable causal analyses in Lamport's happens-before
        sense~\cite{Lampo1978:TimeClocksOrdering},
        from symptoms to root cause.
        Symptoms, whether correctness or performance, materialize at global synchronization boundaries
        where global state becomes meaningfully defined.
        Understanding the root cause of behavior surfacing at these boundaries requires traversing causal
        edges backwards, from symptoms derived from global state to local events capturing their root
        causes.
        A single timestep dataframe captures
        the complete set of events needed to undertake these analyses ---
        those derived from global state surface symptoms, while local events surface the root causes.
        Note that the timestep dataframe abstraction by itself does not guarantee adequate instrumentation
        depth,
        it only provides a generalized windowing construct for correctness and performance analyses.

  \item Control semantics similarly do not create new guarantees, but only preserve pre-existing ones
        in BSP execution.
        Hyperparameters within a BSP workload are strongly consistent across ranks.
        Updates to these hyperparameters are typically coordinated by the lead rank
        via \texttt{MPI\_Bcast()}: a synchronous collective that blocks execution until all ranks receive
        updates.
        However, human operators and other agents attempting to interactively modify these parameters are
        asynchronous entities outside the application domain.
        Any steering inputs that affect execution must be propagated in a
        timestep-consistent manner, a guarantee ORCA provides with \twopc{}.

\end{itemize}




\pplushead{Perturbation Sources and Impact}
All measurement necessarily imposes some perturbation, and is useful insofar as
the induced noise does not obscure the observed signal.
% - On Perturbation - Overall no such thing as observability that does't perturb. more useful thing
% is to ensure that it collects more signal than the noise it induces. To understand ORCA's impact
% on the application, useful to decompose it into three layers -- Instrumentation * ORCA largely
% agnostic. Multiple techniques exist: finstrument-functions, uprobes, kprobes, ring buffers, sync,
% async, Intel PT, CUPTI. Overhead depends on technique. * ORCA's workflow makes better decisions
% achievable, as instrumentation can be disabled if permitted by technique * Overhead in the worst
% case is still unbounded. If an inner loop function is traced, volume can be massive. -- Transport
% Three ways it can bite you. 1. serialization. it is in the hook and can be measured, and arrow is
% super cheap. zero copy is hard but one copy of arrow is the next best thing. 2. jitter from rdma
% pulls. good thing is that this is the best way to retrieve data. esp with QoS. 3. backpressure if
% AGGs backed up. this is a physics bottleneck and not strictly in scope. if you are running out of
% bandwidth either reduce your data collection or provision more AGGs. ORCA provides the design
% knobs. may be reasonable to buffer overflow and drop data, that's a product decision. -- Analytics
% Cost of pushed down kernels. Can be expensive theoretically. Multiple solutions possible. Either
% let users find out the cost of expensive kernels and manually place them at a tier. Or do it
% automatically using the same feedback loop. Kernels are executed synchronously in PTA. Async
% execution is possible for GPU codes, that's an ongoing extension. Sychronous execution cost is
% measurable because it's in the hook. (Overall: there can be two kinds of costs: in the hook or
% outside the hook. jitter from RDMA get is outside the hook, harder to observe. all ops scheduled
% inside the hook are timed byt he infra.)
Perturbation in ORCA arises from either instrumentation, transport, or analytics.
It manifests either as easily measured foreground costs
within the \ocpta{} hook, or
asynchronous work during application execution, which is harder to observe directly.
As its magnitude is highly context-dependent, we now discuss the circumstances under
which each source materially perturbs execution.

\begin{itemize}
  \item \textbf{Instrumentation.
        }
        At its core, ORCA is a streaming analytics framework and agnostic to instrumentation techniques.
        Different instrumentation techniques offer different tradeoffs --- compiled tracepoints
        are cheap but lack runtime flexibility, while eBPF \texttt{uprobes} enable dynamic instrumentation
        but incur a relatively expensive kernel trap per invocation~\cite{Zheng2025:Uprobes}.
        Costs are dependent on the technique used and the functions instrumented,
        and tracing hot inner-loop functions may impose unbounded overhead incurred inline during the execution.
        ORCA's online control capabilities enable
        deactivating expensive probes instantaneously if found to be destabilizing.

  \item \textbf{Transport.
        }
        Transport effects may show as Arrow serialization costs or
        backpressure-driven stalls within the hook, or \texttt{RDMA GET}-induced jitter mid-timestep.
        Arrow serialization is efficient as it involves contiguous buffer \memcpy{}'s, and
        backpressure does not trigger if aggregators are adequately provisioned
        for workload requirements.
        \texttt{RDMA GET} jitter is low as our experiments show, and can be lower
        with modern fabrics' QoS capabilities~\cite{DeSe2020:Slingshot}.

  \item \textbf{Analytics.
        }
        Analytics costs manifest as pushed-down kernels executed on ranks
        within the hook, and are proportional to kernel complexity.
        Our experiments show that these costs are negligible
        for a large class of observability workflows even with highly contended CPU ranks, and the headroom
        only increases with
        GPU-based codes with relatively idle CPU cores.
        However, there is no upper bound on perturbation for arbitrary pushed-down kernels.
        The planner can be augmented
        to reactively adjust placement in response to observed performance impact.

\end{itemize}

Overall, ORCA's architectural and workflow capabilities enable realizing
any given observability task
with costs approaching fundamental hardware limits.
Architecturally, Arrow enables highly efficient columnar
analytics with SIMD kernels and cheap (de)-serialization, RDMA enables low-overhead transport,
while pushdowns minimize data movement required to materialize views.
At a workflow level, steerability enables
dynamically instantiating only the workflows necessary at any given moment.


\pplushead{Architectural Limits of Auxiliary Analytics Tiers}
While elastic, higher TBON tiers typically have lower total I/O bandwidth compared to
the application.
ORCA currently never places sinks on \tiermpi{} ---
while the right default for observability, this may be suboptimal when the higher I/O bandwidth
of the application tier is genuinely necessary, such as during checkpointing.
For observability, offloading sinks minimizes application jitter during the compute phase,
but during synchronous checkpoints, ranks are idle, jitter is not a concern, and maximizing
I/O bandwidth is more important.
ORCA's planner can be extended to be phase-aware to
accommodate streaming analytics use cases beyond observability, and allow placing sinks on
\tiermpi{} when I/O bandwidth becomes a
bottleneck.


% Placing sinks on aggregators is the right default during compute phases because it minimizes
% application jitter by keeping persistence and heavier processing off the MPI hot path. But the AGG
% tier is still auxiliary-sized by construction: its collective I/O bandwidth and compute remain
% much smaller than the aggregate resources of the MPI tier. For observability, this is usually
% acceptable because pushdowns drastically reduce data volume before sink placement becomes the
% bottleneck. For broader in-situ or streaming analytics, however, that assumption can fail. When
% sink volume remains high, the auxiliary tier becomes the limit. This is not a limitation of ORCA's
% architecture so much as of one operating point. In regimes with greater tolerance for
% perturbation, such as checkpointing, it may be preferable to place some sinks or analytics
% directly on MPI and trade additional jitter for much higher aggregate bandwidth. A smarter planner
% can treat MPI and AGG as alternative execution tiers rather than assuming that sinks should always
% be placed on the auxiliary tier.

% \pplushead{When To Transpose} ORCA currently transposes immediately. This requires appending
% elements of a struct to multiple columns. Cache line implications. While in practice we've found
% this to be better than appending a struct sequentially and transposing in bulk, this may cause
% disruption for certain sensitive kernels. When to transpose becomes relevant.

\pplushead{Columnar Transpose Tradeoffs}
Unlike event-based streaming systems, ORCA maintains a columnar structure end-to-end from
generation
to persistence, but this may not be optimal for some high-volume, cache-sensitive functions.
Traces are intrinsically events, which ORCA transposes into
per-column buffers at generation time.
As opposed to \memcpy{}-ing an entire struct
into a \emph{row-major} layout,
this involves small random writes to multiple cache lines to produce a \emph{column-major} layout,
which may perturb some cache-sensitive kernels.
The alternative is a bulk transpose, where events are accumulated in a row-major format,
and transposed in a bulk secondary pass before analytics, at the cost of an additional pass.
In our benchmarks, bulk
transpose offers no measurable benefits and imposes an additional
traversal cost.
However, this effect may be architecture- and
application-dependent.
